\documentclass[a4paper]{article}
\usepackage{fullpage}
\usepackage[color]{vdmlisting}
\usepackage[hidelinks]{hyperref} 
\usepackage{longtable}

\usepackage[portuguese]{babel}
\usepackage[utf8]{inputenc}
\usepackage{indentfirst}
\usepackage{graphicx}
\usepackage{verbatim}
\usepackage{hyperref}
\begin{document}

\title{\Huge\textbf{Green Way System 
}\linebreak\linebreak\linebreak
\Large\textbf{Relatório Final}\linebreak\linebreak
\linebreak\linebreak
\includegraphics[scale=0.1]{feup-logo.png}\linebreak\linebreak
\linebreak\linebreak\linebreak
\linebreak
\Large{Mestrado Integrado em Engenharia Informática e Computação} \linebreak\linebreak
\Large{Métodos Formais em Engenharia de Software}\linebreak\linebreak
\linebreak
}

\author{Ana Rita Ferreira - 201205014/ei12052@fe.up.pt \\ Jorge Teixeira - 201205117/ei12030@fe.up.pt \\\linebreak\linebreak \\
 \\ Faculdade de Engenharia da Universidade do Porto \\ Rua Roberto Frias, s\/n, 4200-465 Porto, Portugal \linebreak\linebreak\linebreak
\linebreak\linebreak\vspace{1cm}}


\maketitle
\thispagestyle{empty}

%************************************************************************************************

\newpage
\tableofcontents
\newpage

\section{Informal system description and list of requirements}
Requirements should include any relevant constraints (regarding safety, etc.). 
Each requirement should have an identifier. 
You may have optional requirements. 
\section{Visual UML model}
A use case model, describing the system actors and 
use cases, with a short description 
of each major use case. 
One or more class diagram(s), describing the structure of the VDM++ model, with a 
short description of each class, plus any other relevant explanations. 
\section{Formal VDM++ model}
VDM++ classes, properly commented. 
Needed data types (e.g., String, Date, etc.) should be modeled with types, values and 
functions. 
Domain entities should be modeled with classes, ins
tance variables and operations. 
You are expected to make adequate usage of the VDM++ types (sets, sequences, maps, etc.) and create a model at a high level of abstraction. 
The model should contain adequate contracts, i.e., 
invariants, preconditions, and -conditions. Post-conditions need only be defined in cases where they are significantly different from the operation or function body (e.g., t
he post-condition of a sqrt(x) operation, which simply states that x = RESULT * RESULT, should be significantly different than the body); for learning purposes, you should define post-conditions for at least two operations. 
During the development of the project, if you foresee that the size of the VDM++ model 
will be less than 5 pages (or 7.5 pages in case of 
groups of 3 students) or more than 10 
pages (or 15 pages in case of groups of 3 students), you should contact your teacher to 
possibly adjust the scope of the system or the modeling approach being followed. 
\section{Model validation (i.e., testing)}
VDM++ test classes, containing adequate and thoroug
h test cases defined by means of 
operations or traces.  
o
Evidences of test results (passed/failed) and test 
coverage. It is sufficient to present the 
system classes mentioned in 4 painted with coverage
 information. Ideally, 100% cov-
erage should be achieved. 
Optionally, figures of examples exercised in the test cases. 
Requirements  traceability  relationship  between  test  cases  and requirements.  Ideally, 100\%  requirements  coverage  should  be  achieved.  It  is  sufficient  to  indicate  in  com-
ments the requirements that are exercised by each test. 
\section{Model verification (i.e., consistency analysis)}
An example of domain verification, i.e., a proof sketch that a pre-condition of an op-
erator, function or operation is not violated. You 
should present the proof obligation 
generated by the tool and your proof sketch.  
An example of invariant verification, i.e., a proof sketch that the body of an operation 
preserves  invariants.  You  should  present  the  proof obligation  generated  by  the  tool 
and your proof sketch. 
\section{Code generation}
You should try to generate Java code from the VDM++ model and try to execute or test 
the generated code. Here you should describe the steps followed and results achieved. 
\section{Conclusions}
Results achieved  
Things that could be improved 
Division of effort and contributions between team members 
\section{References}
\section{Charge}
\begin{vdmpp}[breaklines=true]
class Charge
types
-- TODO Define types here
values
-- TODO Define values here
instance variables
-- TODO Define instance variables here
operations
-- TODO Define operations here
functions
-- TODO Define functiones here
traces
-- TODO Define Combinatorial Test Traces here
end Charge
\end{vdmpp}

\section{Client}
\begin{vdmpp}[breaklines=true]
class Client
types
 public ClientName = seq1 of char;
 public ClientCardNumber = seq1 of nat;
 
values
-- TODO Define values here

instance variables
 private name: ClientName;
 private cardNumber : ClientCardNumber;
 
operations
(*@
\label{Client:14}
@*)
 public Client : ClientName * ClientCardNumber ==> Client
 Client(n, cn) == (
  name := n;
  cardNumber := cn;
  return self
 );
 
(*@
\label{getName:21}
@*)
 pure public getName: () ==> ClientName
 getName() == return name;
 
(*@
\label{getCardNumber:24}
@*)
 pure public getCardNumber: () ==> seq1 of nat
 getCardNumber() == return cardNumber;
 
functions
-- TODO Define functiones here
traces
-- TODO Define Combinatorial Test Traces here
end Client
\end{vdmpp}

\section{Device}
\begin{vdmpp}[breaklines=true]
class Device
types
 public Plate = seq1 of char;
 public ClientCardNumber = seq1 of nat;
values
-- TODO Define values here
instance variables
 private plate:Plate;
 private cardN:ClientCardNumber;
operations
(*@
\label{Client:11}
@*)
(*@
\label{Device:11}
@*)
  public Device : Plate * ClientCardNumber ==> Device
   Device(pt, cn) == (
    plate := pt;
    cardN := cn;
    return self
   );
functions
-- TODO Define functiones here
traces
-- TODO Define Combinatorial Test Traces here
end Device
\end{vdmpp}

\section{Green\_Way}
\begin{vdmpp}[breaklines=true]
class Green_Way
types
-- TODO Define types here
values
-- TODO Define values here
instance variables
-- TODO Define instance variables here
operations
-- TODO Define operations here
functions
-- TODO Define functiones here
traces
-- TODO Define Combinatorial Test Traces here
end Green_Way
\end{vdmpp}

\section{Passage}
\begin{vdmpp}[breaklines=true]
class Passage
types
-- TODO Define types here
values
-- TODO Define values here
instance variables
  public device: Device;
  public provider: [Service_Provider];
  public cost: real;
operations
(*@
\label{Passage:11}
@*)
 public Passage : Device * [Service_Provider] * real ==> Passage
   Passage(dev,sprov,cst) == (
    device := dev ;
   provider := sprov;
   cost := cst;
   return self;
  )
functions
-- TODO Define functiones here
traces
-- TODO Define Combinatorial Test Traces here
end Passage
\end{vdmpp}

\section{Protocol}
\begin{vdmpp}[breaklines=true]
class Protocol
types
-- TODO Define types here
values
-- TODO Define values here
instance variables
-- TODO Define instance variables here
operations
-- TODO Define operations here
functions
-- TODO Define functiones here
traces
-- TODO Define Combinatorial Test Traces here
end Protocol
\end{vdmpp}

\section{Service}
\begin{vdmpp}[breaklines=true]
class Service
types
-- TODO Define types here
values
-- TODO Define values here
instance variables
-- TODO Define instance variables here
operations
-- TODO Define operations here
functions
-- TODO Define functiones here
traces
-- TODO Define Combinatorial Test Traces here
end Service
\end{vdmpp}

\section{Service\_Provider}
\begin{vdmpp}[breaklines=true]
class Service_Provider

instance variables
 private name: seq1 of char;
 private services: set of Service := {};
 
operations
(*@
\label{Service:Provider:8}
@*)
 public Service_Provider : seq1 of char ==> Service_Provider
 Service_Provider(n) == (
  name := n;
  return self
 );
 
(*@
\label{addService:14}
@*)
 public addService : Service ==> ()
 addService(serv) == (
  
    services := services union {serv};
 );
 
(*@
\label{getAllSpots:20}
@*)
 public getAllSpots : () ==> set of Spot
 getAllSpots() == (
    dcl spots: set of Spot := {};
    
    for all service in set services do
     spots := spots union service.getAllSpots();
     
   return spots;
 );
 
 --retorna o custo da passagem
(*@
\label{passa:31}
@*)
 public passa : Client * Spot * Time * [Passage] ==> real
 passa(client, spot, time, last_passage) == (
  
   dcl all_spots : set of Spot := {};
   dcl used_service : Service;
  
    for all service in set services do(
    all_spots := service.getAllSpots();
    if(spot in set all_spots) then 
     used_service := service;
   );
  
   return used_service.passa(client, spot, time, last_passage);
 );
 
end Service_Provider
\end{vdmpp}


\end{document}
